\section{Rangkuman}
\begin{itemize}
    \item Kriptografi adalah ilmu dan seni menjaga keamanan pesan melalui teknik matematika; istilahnya berasal dari \textit{cryptos} (rahasia) dan \textit{graphein} (tulisan).
    \item Empat masalah keamanan komunikasi pada saluran publik adalah kerahasiaan, otentikasi, integritas data, dan nir-penyangkalan.
    \item Sistem kriptografi dimodelkan sebagai quintuple $(\mathcal{E},\mathcal{D},\mathcal{M},\mathcal{K},\mathcal{C})$ yang memetakan plainteks, kunci, dan cipherteks.
    \item Enkripsi mengubah plainteks menjadi cipherteks dengan kunci, sedangkan dekripsi mengembalikannya; keamanan bergantung pada kerahasiaan kunci.
    \item Kriptanalisis adalah upaya memecahkan cipherteks tanpa kunci, dan kriptologi mencakup kriptografi sekaligus kriptanalisis.
    \item Ukuran ruang kunci menentukan ketahanan terhadap \textit{brute-force}: 26 kunci Caesar seketika dipecahkan, sedangkan $26!$ atau $2^{56}$ kunci menuntut waktu yang jauh lebih besar.
    \item Urgensi kriptografi nyata tercermin pada kasus kebocoran data, penyadapan komunikasi, dan pengamanan transaksi elektronik (HTTPS/TLS).
\end{itemize}
